\begin{description}
\item[1.] The trivial solution of the problem is the North Pole ($N$). Starting
  from here we travel 6378\,km to the South ($\overset{\frown}{NA}$) along a meridian
  (e.g.\ Greenwich meridian), then 6378\,km to the East along a parallel
  ($\overset{\frown}{AB}$), and finally 6378\,km to the North ($\overset{\frown}{BN}$) (see
  figure~(a), where the black dashed lines are the Equator and the Greenwich
  meridian).

  For the general solution, mark the path to be taken in one direction with $s$
  and determine what it may be!

  Let us denote the longitude and latitude of the point $A$ by $\lambda_A$ and
  $\varphi_A$, and the point $B$ by $\lambda_B$ and $\varphi_B$.

  The triangle $ANB$ is a spherical one only if the points $A$ and $B$ are on
  the Equator (see figure~(b)), i.e.\ $\varphi_A = \varphi_B = 0^\circ$ and
  $s = R\pi/2$, because the arc $\overset{\frown}{AB}$ is lying along a great circle
  only in this case, otherwise it is a part of a parallel. The point $A$ and
  $B$ can be below the Equator as well (figure~(c)).

  % FIGURE NEEDED: three globe diagrams (a),(b),(c) showing the path N->A->B->N
  % for s = 6378.0 km, 10018.5 km and 10700.0 km; 2019 theory solutions,
  % pages 14-15 (problem 10).

  Next, we measure the angles in radians.

  The central angle of the arc $\overset{\frown}{NA}$ on the meridian fitting to the
  point $A$ ($O$ is the centre of the sphere):
  \[
    NOA\sphericalangle = \frac{\pi}{2} - \varphi,
    \quad\text{where } \varphi = \varphi_A = \varphi_B \tag{10.1}
  \]
  Thus the length $s$ of the arc $\overset{\frown}{NA}$:
  \[
    s = R\left(\frac{\pi}{2} - \varphi\right),
    \quad\text{where } s < R\pi \tag{10.2}
  \]
  The radius of the latitude circle fitting to the points $A$ and $B$:
  \[
    r_{AB} = R\sin\left(\frac{\pi}{2} - \varphi\right) = R\cos\varphi
    \tag{10.3}
  \]
  The central angle of the arc $\overset{\frown}{AB}$ is $\lambda_B - \lambda_A$,
  thus the length of the arc $\overset{\frown}{AB}$:
  \[
    s = r_{AB}\,(\lambda_B - \lambda_A) = R\cos\varphi\,(\lambda_B - \lambda_A)
    \tag{10.4}
  \]
  From the Eq.~(10.2) and Eq.~(10.4):
  \[
    \varphi = \frac{\pi}{2} - \frac{s}{R}
    \quad\text{and}\quad
    \lambda_B - \lambda_A = \frac{s}{R\cos\varphi} \tag{10.5}
  \]
  Substitute the expression of $\varphi$ from the first equation into the
  second equation:
  \[
    \lambda_B - \lambda_A = \frac{s}{R\cos\varphi}
      = \frac{\dfrac{s}{R}}{\cos\left(\dfrac{\pi}{2} - \dfrac{s}{R}\right)}
      = \frac{\dfrac{s}{R}}{\sin\dfrac{s}{R}} \tag{10.6}
  \]
  For a given arc length of $s$ ($s < R\pi$) the latitude $\varphi$ of the
  points $A$ and $B$, and the difference between their longitudes we can
  calculate from the following equations:
  \[
    \varphi = \frac{\pi}{2} - \frac{s}{R}
    \quad\text{and}\quad
    \lambda_B - \lambda_A = \frac{\dfrac{s}{R}}{\sin\dfrac{s}{R}} \tag{10.7}
  \]

  ***** Supplement. It is just for the sake of completeness, not necessary for
  the full marks.

  It is well known, that
  \[
    \lim_{x\to 0}\frac{\sin x}{x} = \lim_{x\to 0}\frac{x}{\sin x} = 1,
    \tag{10.8}
  \]
  thus $\lambda_B - \lambda_A = 1$ in case of $s \to 0$.

  Therefore the difference in longitude has a lower limit (1), but it has no
  upper limit. This means, that the distance between the points $A$ and $B$
  will increase by increasing $s$, the difference between their longitudes will
  approximate $2\pi$ (see figure~(d)), and in case of
  $\lambda_B - \lambda_A = 2\pi$ the points $A$ and $B$ will be the same (see
  figure~(e)). With further increasing of $s$ the difference
  $\lambda_B - \lambda_A$ will be larger than $2\pi$ (see figure~(f)), thus we
  have to stop at that point on the circle where the length of the path along
  it reaches the value of $s$, and then turn to the North.

  % FIGURE NEEDED: three globe diagrams (d),(e),(f) with paths circling the
  % South Pole for s = 17000.0 km, 17206.566 km and 17600.0 km; 2019 theory
  % solutions, page 15 (problem 10).

  The arc length $s_{\lim}$ corresponding to the ``limit case'' of
  $\lambda_B - \lambda_A = 2\pi$ we can derive from the numerical solution of
  the following transcendent equation:
  \[
    \frac{s_{\lim}}{R} - 2\pi\sin\frac{s_{\lim}}{R} = 0
  \]
  Thus find the root of the function
  \[
    f(s) = \frac{s}{R} - 2\pi\sin\frac{s}{R},
  \]
  i.e.\ the value $s_{\lim}$, where $f(s_{\lim}) = 0$. We can solve it by
  iteration using e.g.\ the Newton--Raphson method:
  \[
    s_{k+1} = s_k - \frac{f(s_k)}{f'(s_k)}, \qquad k = 0, 1, 2, \ldots
  \]
  The first derivative of the function $f(s)$:
  \[
    f'(s) = \frac{1}{R} - \frac{2\pi}{R}\cos\frac{s}{R}
  \]
  We know that the parallel in question runs south from the Equator, thus a
  good initial value for the iteration could be $s_0 = R\pi/2 + \Delta s$,
  where $\Delta s = 1000$\,km.

  The result of the iteration and the corresponding latitude with three decimal
  places:
  \[
    s_{\lim} = 17\,206.566\ \mathrm{km}, \quad
    \varphi_{\lim} = -64.573^\circ
  \]
  $\lambda_B - \lambda_A$ is strictly increasing function of $s$ and
  \[
    \lim_{s/R\to 0} \lambda_B - \lambda_A = 1, \quad
    \lim_{s/R\to \pi} \lambda_B - \lambda_A = \infty, \tag{10.9}
  \]
  thus for any values of $0 < s < R\pi$ can be found an appropriate angle
  $\lambda_B - \lambda_A$.

  % FIGURE NEEDED: plot of (lambda_B - lambda_A) = (s/R)/sin(s/R) versus s/R
  % with s_lim/R marked; 2019 theory solutions, page 15 (problem 10).

  ***** End of supplement.

  In conclusion, the solution in this case is the North Pole ($N$).
  \hfill(2\,p)

  For the value of $s = 6378$\,km the solution is an (a)-type curve. The
  coordinates of the turning points are:
  \[
    \lambda_A = 0^\circ, \quad \lambda_B = 68.09^\circ, \quad
    \varphi_A = \varphi_B = 32.704^\circ
  \]
  \hfill(3\,p)

\item[2.] The non-trivial solutions are similar to the (e)-type (``lasso'')
  solution of the previous part, only the starting point is not the North Pole,
  but its location depends on the given value of arc length $s$. The circle of
  the lasso will run south from the (e)-type ``limit circle'', because the
  circumference of parallels running north from it are greater than the
  north--south arc.

  The procedure is the following.

  Knowing $s$, we have to find the parallel around the South Pole ($S$), whose
  circumference is an integer multiple of $s$. \hfill(2\,p)

  If the corresponding latitude is $\varphi_D$, then:
  \[
    2k\pi R\cos\varphi_D = s, \quad k = 1, 2, 3, \ldots \tag{10.10}
  \]
  \hfill(5\,p)

  If the latitude of the sought starting point $C$ is $\varphi_C$, then the
  length $s$ of the arc to be taken along the meridian:
  \[
    s = R\,(\varphi_C - \varphi_D) \tag{10.11}
  \]
  From the Eqs.~(10.10) and (10.11):
  \[
    \varphi_C = \varphi_D + 2k\pi\cos\varphi_D, \quad k = 1, 2, 3, \ldots
    \tag{10.12}
  \]
  Thus:
  \[
    \varphi_D = -\arccos\frac{s}{2k\pi R}, \quad k = 1, 2, 3, \ldots
    \tag{10.13}
  \]
  \hfill(2\,p)

  and
  \[
    \varphi_C = -\arccos\frac{s}{2k\pi R} + \frac{s}{R},
    \quad k = 1, 2, 3, \ldots \tag{10.14}
  \]
  \hfill(2\,p)

  The negative sign denotes that the parallel is in the southern hemisphere.

  So if the starting point is not the North Pole ($N$), there are infinitely
  many solutions.

  For the value of $s = 6378$\,km the first 5 ($k = 1, \ldots, 5$) value for
  $\varphi_D$ in degrees with 3 decimal places:
  \[
    \varphi_{D_1} = -80.842^\circ,\ \varphi_{D_2} = -85.436^\circ,\
    \varphi_{D_3} = -86.959^\circ,\ \varphi_{D_4} = -87.72^\circ,\
    \varphi_{D_5} = -88.176^\circ
  \]
  \hfill(2\,p)

  The first 5 ($k = 1, \ldots, 5$) value for $\varphi_C$ in degrees with 3
  decimal places:
  \[
    \varphi_{C_1} = -23.546^\circ,\ \varphi_{C_2} = -28.14^\circ,\
    \varphi_{C_3} = -29.663^\circ,\ \varphi_{C_4} = -30.424^\circ,\
    \varphi_{C_5} = -30.88^\circ
  \]
  \hfill(2\,p)

  In conclusion, starting to the South from a point of the parallel marked with
  a dashed red line ($C$), we get to a point on a parallel marked with a
  continuous red line ($D$). On this parallel we travel a full circle to the
  East, then we get back to the point $D$, then turn to North and travel back
  to starting point $C$.

  For the further values of $k$ the procedure is very similar, but we have to
  perform two full circles along the continuous green parallel, three along the
  blue, four along the gold, five along the brown and so on.

  Such kind of solution exists for any values of $0 < s < s_{\lim}$.

  % FIGURE NEEDED: globe diagram (g) showing the "lasso" solutions around the
  % South Pole for s = 6378.0 km with the parallels C and D_k marked; 2019
  % theory solutions, page 17 (problem 10).
\end{description}
